\subsection{AgentXploit：黑盒 Agent 的端到端自动化 Red Teaming}

slides 56--62 介绍 AgentXploit。攻击者无法修改用户 query，也不知道内部实现，只能操纵 Agent 会读取的外部环境。框架以 seed attacks 为起点，通过 fuzzing 生成、变异和筛选 payload，在 AgentDojo 与 VWA-adv 上评估有效性、覆盖和迁移，并展示在真实 web agent 的 customer review 中注入内容，诱导 Agent 完成攻击目标。

\lecturefigure{slides-images/slide-056.png}{官方 slide 56：AgentXploit 端到端黑盒 red teaming}
\lecturefigure{slides-images/slide-057.png}{官方 slide 57：AgentXploit motivation 与 threat model}
\lecturefigure{slides-images/slide-058.png}{官方 slide 58：Fuzzing-based methodology 主流程}
\lecturefigure{slides-images/slide-059.png}{官方 slide 59：初始语料、变异与筛选创新}
\lecturefigure{slides-images/slide-060.png}{官方 slide 60：在 AgentDojo 与 VWA-adv 上的评测设置}
\lecturefigure{slides-images/slide-061.png}{官方 slide 61：攻击有效性与迁移评测}
\lecturefigure{slides-images/slide-062.png}{官方 slide 62：真实 Web Agent 的 Customer Review 注入示例}

\begin{knowledgebox}{读图：Fuzzing 为什么适合 Agent attack surface}
Agent 行为由 prompt、页面、memory、tool 和当前状态共同决定，手写少量 payload 很难覆盖组合空间。Fuzzer 通过保留有效 seed、变异表达和依据执行反馈筛选，逐步搜索触发条件。它发现的是具体系统路径，不自动说明所有环境都同样脆弱。
\end{knowledgebox}

\begin{warningbox}{Red Teaming 工具本身也是双用途能力}
自动生成与迁移攻击会降低防守评测成本，也会降低攻击门槛。运行环境应隔离，payload 与成功轨迹受访问控制，测试目标需授权，结果披露遵守协调修复流程；不能把“用于研究”当作无限制发布攻击链的理由。
\end{warningbox}

